\documentclass[11pt,reqno]{amsart}

\usepackage[T1]{fontenc}
\usepackage{lmodern}
\usepackage{amsmath,amssymb,mathtools}
\usepackage{booktabs,array}
\usepackage{enumitem}
\usepackage{microtype}
\usepackage{needspace}
\usepackage[colorlinks=true,linkcolor=blue,citecolor=blue,urlcolor=blue]{hyperref}

\newtheorem{theorem}{Theorem}[section]
\newtheorem{proposition}[theorem]{Proposition}
\newtheorem{lemma}[theorem]{Lemma}
\newtheorem{corollary}[theorem]{Corollary}
\theoremstyle{definition}
\newtheorem{definition}[theorem]{Definition}
\newtheorem{example}[theorem]{Example}
\theoremstyle{remark}
\newtheorem{remark}[theorem]{Remark}

\DeclareMathOperator{\Ker}{Ker}
\newcommand{\F}{\mathbf F}
\newcommand{\N}{\mathbf N}
\newcommand{\Q}{\mathbf Q}
\newcommand{\Z}{\mathbf Z}
\newcommand{\Zp}{\mathbf Z_p}
\newcommand{\Qp}{\mathbf Q_p}
\newcommand{\C}{\mathbf C}

\title[Profinite Mahler recognition]{Profinite Recognition of Mahler Equations\\
and D-Finite Witt Escape}
\author{KEIO}
\date{September 2026}
\hypersetup{
  pdftitle={Profinite Recognition of Mahler Equations and D-Finite Witt Escape},
  pdfauthor={KEIO}
}

\begin{document}

\begin{abstract}
Fix an integer base $k\geq2$ and a power series $F(t)$ with rational
coefficients.  We prove
that Mahler relations in a fixed order--degree box are recognized by
arbitrarily accurate prime-power reductions, even when the modular
relations are mutually incompatible.  At a fixed prime the exact maximal
false-positive depth is the stabilized largest Smith exponent of a
sufficiently large full-rank coefficient truncation.  Across primes, one
nonzero minor gives a logarithmic precision budget, and divergence of the
total local precision is equivalent to an exact relation.

We also prove a nonlinear recognition theorem: if, to unbounded $p$-adic
depth, a bounded-degree multiplier makes the full Cartier matrix have
bounded determinantal rank, then $F$ is a characteristic-zero Mahler
function.  Combining these mechanisms with classical theorems gives three
rigidity consequences.  Persistence in two multiplicatively independent
bases is equivalent to rationality; every nonrational D-finite series
escapes each fixed Witt--Mahler box; and its reduced $p$-kernel cardinalities
tend to infinity.  When the reductions are automatic, the last statement
is an automaton-state lower conclusion.

An adelic one-pole hypothesis gives a further unconditional consequence:
the reciprocal of the unique pole in the complex open unit disk is a Pisot
or Salem number.  Finally, sharp boundary examples separate fixed-base
depth from changing-characteristic phenomena.  The central-binomial series
has full own-$p$ Cartier rank one modulo $p$, while the least coefficient
degree of an endpoint-nonzero own-$p$ relation is exactly $(p-1)/2$.  We
also exhibit a series rational modulo every $p^m$ but non-Mahler in every
integer base.  No effective escape rate is claimed.
\end{abstract}

\maketitle

\section{Introduction}

A nonzero formal power series $F$ is called $k$-Mahler if it satisfies an
identity
\begin{equation}\label{eq:mahler}
 \sum_{i=0}^{q}P_i(t)F(t^{k^i})=0,
 \qquad P_i(t)\in\Q[t],\qquad P_0P_q\ne0.
\end{equation}
We regard the zero series as Mahler in every base.  A relation initially
missing its $F(t)$ endpoint can be normalized to this form; the elementary
residue-section argument is recorded in Lemma~\ref{lem:endpoint} below.
Reduction modulo $p^m$ often produces identities of this shape long before
one knows whether an identity exists in characteristic zero.  The main
question is deliberately uniform: what follows if the base, order and
coefficient degrees remain bounded while arithmetic precision tends to
infinity, vertically at one prime, horizontally through many primes, or by
a mixture of the two?

The answer is exact.  Compactness gives a $p$-adic relation without requiring
the modular coefficient vectors to be compatible, and a finite-minor rank
argument descends it to $\Q$.  Smith normal form then identifies the last
possible false-positive depth.  This linear mechanism has a nonlinear
counterpart in which a bounded polynomial ``blocks'' the poles before a
uniform Cartier-rank constraint is imposed.

We use \emph{Witt depth} only as shorthand for passage through the coefficient
rings $\Z/p^m\Z$.  No operation from the formal theory of Witt vectors is
used.

One consequence no longer requires D-finiteness.  Two exact Mahler
relations in multiplicatively independent bases force rationality by the
theorem of Adamczewski and Bell~\cite{abtwobase}; our recognition theorem
turns this into an if-and-only-if criterion stated solely in terms of
prime-power persistence.  The principal escape consequences concern
D-finite series.  B\'ezivin's theorem,
in the formulation of Bell, Coons and Rowland~\cite{bcr}, says that a series
which is both D-finite and $k$-Mahler is rational.  Hence modular Mahler
relations for a nonrational D-finite series cannot remain in a fixed box
through all Witt depths.  This is a complexity escape statement, not a
quantitative lower bound.

Our results are organized as follows.  Section~2 gives fixed-box profinite
recognition and the endpoint normalization needed by later applications.
Section~3 computes exact local precision, proves mixed-prime recognition,
and derives the two-base rationality criterion.  Section~4 proves the
multiplier--Cartier recognition theorem.  Sections~5 and~6 derive D-finite
and automata consequences.  Section~7 treats the changing-base obstruction,
Section~8 gives sharpness examples, and Section~9 proves the adelic pole
reversal.

The arguments combine elementary compactness and determinantal algebra with
classical inputs of Becker~\cite{becker}, B\'ezivin~\cite{bezivin},
Adamczewski--Bell~\cite{abtwobase}, Adamczewski--Bell--Smertnig~\cite{abs},
and Rowland--Yassawi~\cite{profinite,rowlandyassawi}.  Our contribution is
the exact fixed-box threshold, the incompatible-blocker recognition
principle, and the escape and profinite reformulations obtained by assembling
these tools.  The two-base rationality implication and the Mahler-denominator
input in Section~9 are explicitly credited corollaries of the cited
theorems.  We do not claim a solution of Christol's conjecture, a general
diagonal conjecture, or the Salem--Boyd closure problem.

\section{Fixed-box profinite recognition}

Throughout, $k\geq2$ and $p$ is prime.  Let
\[
 F(t)=\sum_{n\geq0}a_nt^n\in\Q[[t]]\cap\Zp[[t]].
\]
For $q\geq1$ and $h\geq0$ define the finite family
\begin{equation}\label{eq:box-family}
 \mathcal V^{(k)}_{q,h}(F)=
 \{t^jF(t^{k^i}):0\leq i\leq q,\ 0\leq j\leq h\}.
\end{equation}
When the base is fixed we write simply $\mathcal V_{q,h}(F)$.
It has $N=(q+1)(h+1)$ elements.  A coefficient vector in
$(\Z/p^m\Z)^N$ is called \emph{primitive} if at least one coordinate is a
unit modulo $p$.

The next elementary point prevents a hidden endpoint gap when a linear
dependence is converted into a standard Mahler equation.

\begin{lemma}[Endpoint normalization]\label{lem:endpoint}
Let $K$ be a field, $F\in K[[t]]$ be nonzero, and suppose
\begin{equation}\label{eq:shifted-mahler}
 \sum_{i=s}^{q}P_i(t)F(t^{k^i})=0,
 \qquad P_s\ne0,
\end{equation}
with $P_i\in K[t]$.  Then $F$ satisfies a $k$-Mahler equation whose
coefficient of $F(t)$ is nonzero.
\end{lemma}

\begin{proof}
Write, uniquely,
\[
 P_i(t)=\sum_{r=0}^{k^s-1}t^rP_{i,r}(t^{k^s}).
\]
The residue classes of exponents modulo $k^s$ in
\eqref{eq:shifted-mahler} do not interact.  On putting $u=t^{k^s}$, each
residue class therefore gives
\[
 \sum_{i=s}^{q}P_{i,r}(u)F(u^{k^{i-s}})=0.
\]
Choose $r$ with $P_{s,r}\ne0$ and delete any zero terms at the far end.
The resulting relation begins with a nonzero multiple of $F(u)$.  It must
also contain a later term, since otherwise $P_{s,r}F=0$, contrary to the
hypotheses.
\end{proof}

\begin{theorem}[Vertical profinite recognition]\label{thm:vertical}
For fixed $k,q,h$, the following are equivalent.
\begin{enumerate}[label=\textup{(\roman*)}]
\item The family $\mathcal V_{q,h}(F)$ is linearly dependent over $\Q$.
\item For every $m\geq1$ there is a primitive coefficientwise relation among
      the family modulo $p^m$.
\item Such a primitive relation exists for arbitrarily large $m$.
\end{enumerate}
No compatibility between the relations at different depths is required.
\end{theorem}

\begin{proof}
The implication (i)$\Rightarrow$(ii) follows after multiplying a rational
relation by a scalar that makes its coefficients $p$-integral and primitive.
The implication (ii)$\Rightarrow$(iii) is immediate.

For (iii)$\Rightarrow$(i), lift the modular coefficient vectors to the
compact $p$-adic unit shell
\[
 S=\{c\in\Zp^N:\max_j|c_j|_p=1\}.
\]
Along depths tending to infinity, take a convergent subsequence with limit
$c\in S$.  Each coefficient of the formal series
\[
 \sum_{i=0}^{q}\sum_{j=0}^{h}c_{i,j}t^jF(t^{k^i})
\]
is a continuous $\Zp$-linear function of $c$.  Its values along the
subsequence lie in ideals $p^m\Zp$ with $m\to\infty$, so its value at the
limit is zero.  This gives a nonzero dependence over $\Qp$.

The coefficient matrix has rational entries.  If its columns were independent
over $\Q$, some finite rational $N$-by-$N$ minor would be nonzero, and it
would remain nonzero over $\Qp$.  Thus the dependence descends to $\Q$.
\end{proof}

For $F\ne0$, every exact dependence furnished by
Theorem~\ref{thm:vertical} becomes a standard $k$-Mahler equation by
Lemma~\ref{lem:endpoint}.  The zero series is already Mahler by convention.

\begin{remark}\label{rem:uniformity}
Fixed base, fixed $q,h$, and primitivity are essential.  Allowing the box to
grow with $m$, or allowing all coefficients of a relation to be divisible by
$p$, destroys the conclusion.  Section~7 explains a different phenomenon
when the base itself changes with $p$.
\end{remark}

\section{Exact precision and finite certificates}

Let $M=M_{q,h}(F)$ be the infinite coefficient matrix whose $N$ columns are
the series in \eqref{eq:box-family}.  First suppose that $M$ has column rank
$N$.  For a finite row set $R$ on which $M_R$ has rank $N$, view $M_R$ over
$\Zp$ and write its Smith invariant factors as
\[
 d_1(R)\mid\cdots\mid d_N(R),\qquad
 \alpha_R=v_p(d_N(R)).
\]

Set
\begin{equation}\label{eq:local-depth-definition}
 \epsilon_p(M)=\sup\bigl(\{m\geq1:Mx\equiv0\pmod{p^m}
 \text{ has a primitive solution}\}\cup\{0\}\bigr).
\end{equation}

\begin{proposition}[Exact local depth]\label{prop:smith}
The exact maximal primitive depth of the full infinite system is
\begin{equation}\label{eq:epsilon}
 \epsilon_p(M)=\min_R\alpha_R,
\end{equation}
where $R$ ranges over finite full-rank row sets.  Along any increasing
full-rank exhaustion of the rows, the corresponding exponents decrease and
eventually stabilize at $\epsilon_p(M)$.
\end{proposition}

\begin{proof}
Unimodular row and column operations preserve primitivity.  For a finite
matrix with Smith diagonal $d_1,\ldots,d_N$, a primitive kernel vector modulo
$p^m$ exists exactly when $m\leq v_p(d_N)$.  This proves the upper bound in
\eqref{eq:epsilon}.

Conversely, suppose $m\leq\alpha_R$ for every finite full-rank $R$.  Every
finite set of row equations can be enlarged to a full-rank one and hence has
a primitive solution modulo $p^m$.  The finite set of primitive vectors in
$(\Z/p^m\Z)^N$ therefore has the finite-intersection property for these row
constraints, so one vector solves all rows.  Adding rows can only decrease
the exponent.  A nonincreasing sequence of nonnegative integers stabilizes,
which proves the final assertion.
\end{proof}

If $M$ has rank less than $N$, there is an exact relation and the available
depth is infinite.  Proposition~\ref{prop:smith} concerns precisely the
full-rank case.

The same argument provides a global obstruction.  Let $\Delta_0\neq0$ be a
fixed rational $N$-by-$N$ minor of $M$.  Write it in lowest terms.

\begin{proposition}[Single-minor precision budget]\label{prop:budget}
Let $\mathcal P$ be a finite set of primes at which all entries used in
$\Delta_0$ are integral.  If a primitive relation exists modulo $p^{m_p}$
for every $p\in\mathcal P$, then
\begin{equation}\label{eq:budget}
 \sum_{p\in\mathcal P}m_p\log p
 \leq\log|\operatorname{num}(\Delta_0)|.
\end{equation}
In particular, a nonzero minor is an explicit certificate that relations in
this box cannot persist past its local valuations.
\end{proposition}

\begin{proof}
If $Bx\equiv0\pmod{p^{m_p}}$ for the square minor matrix $B$ and $x$ is
primitive, the adjugate identity implies
$v_p(\det B)\geq m_p$.  Hence $p^{m_p}$ divides the numerator of
$\Delta_0$.  Multiplication over $p\in\mathcal P$ and taking logarithms gives
\eqref{eq:budget}.
\end{proof}

There is also an integral finite-truncation version.  For a finite full-rank
$M_R$, let $D_R$ clear its denominators and let $s_N(R)$ be the largest
integer Smith invariant of $A_R=D_RM_R$.  At primes $p\nmid D_R$,
\begin{equation}\label{eq:finite-smith-budget}
 m_p\leq v_p(s_N(R)),\qquad
 \sum_{p\in\mathcal P}m_p\log p\leq\log|s_N(R)|.
\end{equation}
We emphasize the direction: a nonzero minor or Smith invariant gives a
\emph{negative} certificate against deeper lifting.  Without an independent
row cutoff it is not a finite positive certificate for the existence of an
exact relation.

\subsection{Vertical, horizontal and mixed precision}

We now reset the ambient arithmetic.  For $F=\sum a_nt^n\in\Q[[t]]$, call
a prime $p$ \emph{good for $F$} if $a_n\in\Zp$ for every $n$.  At a good
prime, let $\epsilon_p^{(b;q,h)}(F)\in\N\cup\{0,+\infty\}$ be the supremum
of $0$ and all depths $m\geq1$ at which $\mathcal V^{(b)}_{q,h}(F)$ has a
primitive coefficientwise relation modulo $p^m$.

\begin{theorem}[Mixed vertical--horizontal recognition]
\label{thm:mixed-profile}
Let $b\geq2$, $q\geq1$, $h\geq0$, and let $\mathcal P$ be a nonempty set
of primes good for $F\in\Q[[t]]$.  Then
\begin{equation}\label{eq:mixed-equivalence}
 \mathcal V^{(b)}_{q,h}(F)\text{ is dependent over }\Q
 \quad\Longleftrightarrow\quad
 \sum_{p\in\mathcal P}\epsilon_p^{(b;q,h)}(F)\log p=+\infty.
\end{equation}
If the family is independent and $\Delta_0\ne0$ is any rational maximal
minor of its coefficient matrix, then more precisely
\begin{equation}\label{eq:mixed-bound}
 \sum_{p\in\mathcal P}\epsilon_p^{(b;q,h)}(F)\log p
 \leq\log|\operatorname{num}(\Delta_0)|.
\end{equation}
A summand with value $+\infty$ makes the sum infinite.
For an infinite prime set, the sum means the supremum of its finite
sub-sums.
\end{theorem}

\begin{proof}
An exact rational relation can be cleared to a primitive integer vector.  It
is primitive at every prime and gives relations through every depth, proving
the forward implication.  Conversely, suppose the family is independent
and choose $\Delta_0$.  For every finite
$\mathcal P_0\subset\{p\in\mathcal P:\epsilon_p^{(b;q,h)}(F)\geq1\}$,
Proposition~\ref{prop:budget}, applied at the maximal finite available
depths, gives
\[
 \sum_{p\in\mathcal P_0}\epsilon_p^{(b;q,h)}(F)\log p
 \leq\log|\operatorname{num}(\Delta_0)|.
\]
If one local depth is infinite, apply the same argument at arbitrarily large
finite depths.  Taking the supremum over finite $\mathcal P_0$ proves both
\eqref{eq:mixed-bound} and the reverse implication.
\end{proof}

Thus one fixed prime of unbounded depth is the \emph{vertical} case;
infinitely many good primes of positive depth give the \emph{horizontal}
case; and every divergent weighted mixture is allowed.  The relations may
vary with both prime and depth, but the base and the box may not vary.

\begin{theorem}[Two-base profinite rationality]\label{thm:two-base}
Let $k,\ell\geq2$ be multiplicatively independent and let
$F\in\Q[[t]]$.  The following are equivalent.
\begin{enumerate}[label=\textup{(\roman*)}]
\item $F\in\Q(t)$.
\item For each $b\in\{k,\ell\}$ there are fixed $q_b\geq1$, $h_b\geq0$
and a nonempty set $\mathcal P_b$ of primes good for $F$ such that
\[
 \sum_{p\in\mathcal P_b}
 \epsilon_p^{(b;q_b,h_b)}(F)\log p=+\infty.
\]
\end{enumerate}
The two prime sets, depths and modular relations need not agree.
\end{theorem}

\begin{proof}
Assume (ii).  Theorem~\ref{thm:mixed-profile}, followed by
Lemma~\ref{lem:endpoint} when $F\ne0$, makes $F$ both $k$-Mahler and
$\ell$-Mahler.  The theorem of Adamczewski and Bell
\cite[Theorem~1.3]{abtwobase} gives $F\in\Q(t)$.

Conversely, let $F=A/B\ne0$ with $A,B\in\Z[t]$ and $B(0)\ne0$.  For every
$b\geq2$,
\begin{equation}\label{eq:rational-mahler}
 A(t^b)B(t)F(t)-A(t)B(t^b)F(t^b)=0.
\end{equation}
After division by the joint integer content, its coefficient vector is
primitive at every prime.  Every $p\nmid B(0)$ is good for $F$, and
\eqref{eq:rational-mahler} persists through every depth at any such prime.
This proves (ii).  For $F=0$, use the relation $F(t)=0$.
\end{proof}

Theorem~\ref{thm:two-base} is a profinite reformulation obtained by
combining Theorem~\ref{thm:mixed-profile} with the characteristic-zero
two-base theorem of Adamczewski and Bell; no new proof of their rationality
theorem is claimed here.

\section{Polynomial blockers and Cartier rank}

For $0\leq r<k$ define the Cartier section
\[
 \Lambda_{k,r}\!\left(\sum_{n\geq0}b_nt^n\right)
   =\sum_{n\geq0}b_{kn+r}t^n.
\]
For a word $w$ in the digits $0,\ldots,k-1$, let $\Lambda_w$ be the
corresponding iterate.  The full Cartier matrix of a series $G$ is
\begin{equation}\label{eq:cartier-matrix}
 \mathcal C_k(G)=\bigl([t^n]\Lambda_wG\bigr)_{n\geq0,\ w}.
\end{equation}
Its column rank is the dimension of the linear span of the $k$-kernel of
$G$.

\begin{proposition}[Raw Cartier rank cannot preserve an inner pole]
\label{prop:raw-radius}
If $G(t)=\sum_{n\geq0}b_nt^n\in\C[[t]]$ has finite-dimensional full
$k$-Cartier span, then $|b_n|=O(n^c)$ for some $c\geq0$.  In particular,
the complex radius of convergence of $G$ is at least $1$.
\end{proposition}

\begin{proof}
On the finite-dimensional span of the Cartier sections, the $k$ section
operators are represented by finitely many matrices.  Reading the base-$k$
digits of $n$ expresses $b_n$ as a fixed linear functional applied to a
product of $O(\log n)$ of these matrices.  A common operator-norm bound
therefore gives $|b_n|\leq C A^{O(\log n)}=O(n^c)$.
\end{proof}

Thus finite raw Cartier rank cannot hold for a series with a pole strictly
inside the complex unit disk.  The multiplier below is not cosmetic: it may
cancel finitely many poles before rank is measured.

\begin{definition}
A polynomial $Q(t)=\sum_{j=0}^{M}q_jt^j$ over $\Z/p^m\Z$ is a
\emph{primitive $(M,D)$ blocker for $F$ at depth $m$} if some $q_j$ is a
unit modulo $p$ and every $(D+1)$-minor of
$\mathcal C_k(QF)$ vanishes modulo $p^m$.
\end{definition}

\begin{theorem}[Nonlinear profinite recognition]\label{thm:blocker}
Let $p$ be prime, let $F\in\Q[[t]]\cap\Zp[[t]]$, and fix
$k\geq2$ and $M,D\geq0$.  If primitive $(M,D)$ blockers exist for
arbitrarily large $m$, then $F$ is $k$-Mahler over $\Q$.
\end{theorem}

\begin{proof}
Lift blocker coefficients to the compact unit shell in $\Zp^{M+1}$.  For
each $m$, let $X_m$ be the closed subset on which every indicated minor lies
in $p^m\Zp$.  The hypothesis makes $X_m$ nonempty and
$X_{m+1}\subseteq X_m$.  Compactness gives a point
$Q\in\bigcap_mX_m$.  It is nonzero, has degree at most $M$, and every
$(D+1)$-minor of $\mathcal C_k(QF)$ vanishes exactly.  Thus
$G=QF$ has a finite-dimensional $k$-kernel over $\Qp$.

Finite Cartier span means that $G$ is $k$-regular.  Becker's
theorem~\cite[Theorem~1]{becker} then gives a $k$-Mahler equation over
$\Qp(t)$.  Substituting $G(t^{k^i})=Q(t^{k^i})F(t^{k^i})$ and clearing
denominators gives some finite relation among the iterates of $F$ over
$\Qp[t]$.  The corresponding finite coefficient matrix has rational entries,
because $F\in\Q[[t]]$.  The same finite-minor rank descent used in
Theorem~\ref{thm:vertical} produces a relation over $\Q[t]$; if $F\ne0$,
Lemma~\ref{lem:endpoint} puts it in standard Mahler form.
\end{proof}

\begin{remark}
The limiting blocker need not have rational coefficients.  The proof gives
existence of a characteristic-zero Mahler equation but, without an effective
form of the Becker and elimination steps, gives no bound for its eventual
order or coefficient degree in terms of $k,M,D$.
\end{remark}

The nonlinear system also has a horizontal and mixed-prime obstruction,
although its certificate is a projective Nullstellensatz identity rather
than one determinant.  At a good prime, let
$\eta_p^{(b;M,D)}(F)\in\N\cup\{0,+\infty\}$ be the supremum of $0$ and all
depths $m\geq1$ at which a primitive $(M,D)$ blocker in base $b$ exists
modulo $p^m$.

\begin{theorem}[Mixed-prime blocker recognition]\label{thm:mixed-blocker}
Fix $b\geq2$ and $M,D\geq0$.  If $F\in\Q[[t]]$ is not $b$-Mahler, then
there is a nonzero integer $A=A(F;b,M,D)$ such that, for every set
$\mathcal P$ of good primes,
\begin{equation}\label{eq:blocker-budget}
 \sum_{p\in\mathcal P}\eta_p^{(b;M,D)}(F)\log p\leq\log|A|.
\end{equation}
Consequently, divergence of the sum forces $F$ to be $b$-Mahler.
\end{theorem}

\begin{proof}
Regard every $(D+1)$-minor of $\mathcal C_b(QF)$ as a homogeneous
polynomial $f_\nu(q_0,\ldots,q_M)$ over $\Q$.  If all these polynomials had
a common nonzero zero over $\overline\Q$, the corresponding algebraic
polynomial $Q$ would make $QF$ have finite Cartier rank.  Becker's theorem,
substitution, rational-column rank descent and
Lemma~\ref{lem:endpoint} would then make $F$ $b$-Mahler.  Under the stated
hypothesis the projective common zero set is therefore empty.

By Hilbert's basis theorem, finitely many of the $f_\nu$ generate the same
homogeneous ideal.  The projective Nullstellensatz puts a power of each
$q_j$ in its scalar extension to $\overline\Q$; faithful flatness descends
these ideal-membership certificates to $\Q$.  After clearing denominators,
there are an integer $A\ne0$, an exponent $L$, and integer polynomials
$H_{j,\nu}$ such that
\begin{equation}\label{eq:nullstellensatz-certificate}
 A q_j^L=\sum_\nu H_{j,\nu}(q)g_\nu(q)
 \qquad(0\leq j\leq M),
\end{equation}
where each $g_\nu$ is an integer multiple of the corresponding minor.
For a primitive blocker modulo $p^m$, choose a lift
$\widetilde q\in\Zp^{M+1}$.  Goodness of $p$ makes the selected minor
polynomials $p$-integral, some $\widetilde q_j$ is a unit, and all
$g_\nu(\widetilde q)$ vanish modulo $p^m$; hence $p^m\mid A$.  Thus
$\eta_p^{(b;M,D)}(F)\leq v_p(A)$ at every good prime.  Summing over finite
prime sets and taking the supremum proves \eqref{eq:blocker-budget}.
\end{proof}

The algebraic blocker arising in the first paragraph need not have rational
coefficients.  What descends to $\Q$ is the resulting Mahler equation for
the rational-coefficient series $F$.

\begin{corollary}[Two-base blocker criterion]\label{cor:two-base-blocker}
Let $k,\ell\geq2$ be multiplicatively independent and
$F\in\Q[[t]]$.  Then $F$ is rational if and only if, for each
$b\in\{k,\ell\}$, there are fixed $M_b,D_b$, a nonempty set $\mathcal P_b$
of good primes, and
\[
 \sum_{p\in\mathcal P_b}\eta_p^{(b;M_b,D_b)}(F)\log p=+\infty.
\]
\end{corollary}

\begin{proof}
Assume the blocker condition.  Theorem~\ref{thm:mixed-blocker} makes $F$
Mahler in both bases, so Adamczewski--Bell gives $F\in\Q(t)$.  Conversely,
write $F=A/B$ with $A,B\in\Z[t]$ and $B(0)\ne0$.  At every
$p\nmid B(0)$, the primitive multiplier $Q=B$ makes $QF=A$ a polynomial;
all its Cartier sections lie in a space of dimension at most $\deg A+1$.
Thus fixed bounds work at every depth.  The zero series is immediate with
$Q=1$.
\end{proof}

\begin{corollary}[Archimedean obstruction to profinite persistence]
\label{cor:height-obstruction}
Let $F(t)=\sum a_nt^n\in\Q[[t]]$ and let $h(a_n)$ denote logarithmic Weil
height.  If
\[
 h(a_n)=o(n)\qquad\text{but}\qquad h(a_n)\not=O((\log n)^2),
\]
then, in every fixed base, the total precision profile of every fixed
linear box and every fixed blocker box is finite; explicitly, the sums of
$\epsilon_p\log p$ and $\eta_p\log p$ over all good primes are finite.
\end{corollary}

\begin{proof}
Infinite total precision in a linear box would make $F$ Mahler by
Theorem~\ref{thm:mixed-profile}; in a blocker box it would do so by
Theorem~\ref{thm:mixed-blocker}.  The height-gap theorem of Adamczewski,
Bell and Smertnig~\cite[Theorem~6.1]{abs} says that a rational-coefficient
Mahler series of sublinear coefficient height has height
$O((\log n)^2)$, a contradiction.
\end{proof}

The finer height ladder in~\cite{abs} gives analogous obstructions at the
successive thresholds $\log^2 n$, $\log n$, $\log\log n$ and bounded
height.  These conclusions do not require D-finiteness.

\section{D-finite fixed-box escape}

We now assume that $F$ is D-finite over $\Q(t)$.  Recall that a D-finite
$k$-Mahler series in characteristic zero is rational~\cite{bcr}.

\begin{theorem}[D-finite Witt--Mahler dichotomy]\label{thm:dfinite}
Let $F\in\Q[[t]]\cap\Zp[[t]]$ be D-finite and fix $k\geq2$.  Then the
following are equivalent.
\begin{enumerate}[label=\textup{(\roman*)}]
\item $F$ is rational.
\item Some fixed order--degree box admits primitive Mahler relations modulo
      $p^m$ for arbitrarily large $m$.
\item Some fixed pair $(M,D)$ admits primitive polynomial blockers modulo
      $p^m$ for arbitrarily large $m$.
\end{enumerate}
\end{theorem}

\begin{proof}
The implications (ii)$\Rightarrow$(i) and (iii)$\Rightarrow$(i) follow from
Theorems~\ref{thm:vertical} and~\ref{thm:blocker}, respectively, from
Lemma~\ref{lem:endpoint} when needed, and then from the
D-finite--Mahler rationality theorem.

For the converse, the case $F=0$ is immediate: take the relation $F(t)=0$
and the blocker $Q=1$.  If $F=A/B\ne0$ is rational, then
after a rational rescaling we may take $A,B\in\Zp[t]$ with $B$ primitive.
Then
\[
 A(t^k)B(t)F(t)-A(t)B(t^k)F(t^k)=0,
\]
which gives (ii) at every depth after primitive normalization.  Taking
$Q=B$ makes $QF=A$ a polynomial.  Its iterated Cartier sections span a
finite-dimensional space, so a fixed $(M,D)$ gives (iii).
\end{proof}

Define the fixed-base linear complexity at depth $m$ by
\begin{equation}\label{eq:mu}
 \mu_{k,p}(F;m)=
 \min\{(q+1)(h+1):\mathcal V_{q,h}(F)
 \text{ has a primitive relation mod }p^m\},
\end{equation}
with value $+\infty$ when there is none.

\begin{corollary}[Cofinal escape]\label{cor:escape}
If $F\in\Q[[t]]\cap\Zp[[t]]$ is D-finite and nonrational, then
\[
 \mu_{k,p}(F;m)\longrightarrow+\infty.
\]
Moreover, for every fixed $(M,D)$ there is a depth after which no primitive
$(M,D)$ blocker exists.
\end{corollary}

\begin{proof}
Relations at depth $m+1$ reduce to primitive relations at depth $m$, so the
minimum in \eqref{eq:mu} is nondecreasing.  If it remained bounded along
unbounded depths, one of finitely many boxes would recur arbitrarily deeply,
contradicting Theorem~\ref{thm:dfinite}.  The blocker assertion is the same
contrapositive.
\end{proof}

This corollary asserts escape from every fixed box; it supplies no rate and
no lower bound for the order of a P-recursive recurrence.

\section{Automatic state complexity modulo \texorpdfstring{$p^m$}{prime powers}}

For a sequence $a=(a_n)_{n\geq0}$ with values modulo $p^m$, its $p$-kernel
is
\[
 \Ker_p(a)=\{(a_{p^en+r})_{n\geq0}:e\geq0,\ 0\leq r<p^e\}.
\]
When this set is finite, its cardinality is the state count of the normalized
minimal least-significant-digit-first automaton unaffected by leading
zeros~\cite{eilenberg,rowlandyassawi}.

\begin{theorem}[$p$-kernel escape]\label{thm:states}
Let $F(t)=\sum a_nt^n\in\Q[[t]]\cap\Zp[[t]]$ be D-finite and nonrational.
Put, in the extended positive integers,
\[
 N_m=|\Ker_p(a\bmod p^m)|\in\N\cup\{+\infty\}.
\]
Then $N_m\to\infty$.
\end{theorem}

\begin{proof}
Reduction modulo $p^m$ maps the kernel at depth $m+1$ onto the kernel at
depth $m$, so $N_m$ is nondecreasing.  If $N_m\leq B$ for all $m$, the exact
$p$-kernel of the $\Zp$-valued sequence has at most $B$ elements.  Indeed,
$B+1$ distinct exact kernel sequences can be separated at finitely many
coefficient positions; at a sufficiently large depth all those differences
remain visible, contradicting $N_m\leq B$.

Every exact kernel sequence still has rational entries.  If the exact
$p$-kernel is finite, its $\Q$-span is finite-dimensional and stable under
all Cartier sections.  Thus the series is $p$-regular over $\Q$, and
Becker's theorem gives a $p$-Mahler equation over $\Q$.
Lemma~\ref{lem:endpoint} puts it in standard form, and
D-finite--Mahler rationality contradicts the hypothesis.
\end{proof}

No automaticity assumption is needed for the extended-cardinality
conclusion.  To read it as an automaton statement one does need every
$N_m$ to be finite.  This is available for algebraic $p$-adic series under
the prime-power automaticity results summarized in~\cite{profinite}, and
for rational diagonals under the integral hypotheses of
Rowland--Yassawi~\cite{diagonals,rowlandyassawi}.  It is not known for
arbitrary D-finite integer series.

\begin{corollary}[Automatic state escape]\label{cor:algebraic}
Under the hypotheses of Theorem~\ref{thm:states}, assume additionally that
$a\bmod p^m$ is $p$-automatic for every $m$.  Then the normalized minimal
automaton state counts tend to infinity.  In particular this applies to
nonrational algebraic series and to nonrational D-finite rational diagonals
whenever the cited all-depth automaticity hypotheses hold.
\end{corollary}

For a series already known to be integral and algebraic, the conclusion is
especially transparent: finite automata exist at each finite Witt level,
but no uniform finite state space can survive the inverse tower.  The result
is complementary to the explicit upper bounds in~\cite{rowlandyassawi}.

\section{Changing to the residue-characteristic base}

Fixed-base uniformity cannot be replaced by choosing base $p$ at prime $p$.
Nevertheless bounded coefficient degree across infinitely many primes is so
restrictive that it forces a stronger collapse.

\begin{proposition}[Expanding-prefix recurrence lifting]
\label{prop:prefix-lifting}
Let $F=\sum a_nt^n\in\Q[[t]]$.  Suppose there are an integer $h\geq0$ and
infinitely many good primes $p$ for which a nonzero vector
$(c_{p,0},\ldots,c_{p,h})\in\F_p^{h+1}$ satisfies
\begin{equation}\label{eq:prefix-recurrence}
 \sum_{j=0}^{h}c_{p,j}\,\overline a_{n-j}=0
 \qquad(h<n<p).
\end{equation}
Then $F$ is rational.
\end{proposition}

\begin{proof}
Form the infinite rational Toeplitz matrix whose row indexed by $n>h$ is
\[
 (a_n,a_{n-1},\ldots,a_{n-h}).
\]
If its columns had full rank, a
fixed nonzero rational minor would remain nonzero modulo all but finitely
many primes.  For every sufficiently large prime in the assumed set, all
rows of that minor occur in the range of \eqref{eq:prefix-recurrence}, a
contradiction.  Thus the columns are dependent over $\Q$, giving an exact
constant-coefficient recurrence and hence a rational generating function.
\end{proof}

For a good prime $p$, let $L_p(F)$ be the least
$h\in\{0,\ldots,p-2\}$ for which some
$0\ne(c_0,\ldots,c_h)\in\F_p^{h+1}$ satisfies
\[
 \sum_{j=0}^{h}c_j\overline a_{n-j}=0\qquad(h<n<p).
\]
Set $L_p(F)=p-1$ if no such $h$ exists.
Proposition~\ref{prop:prefix-lifting} immediately gives
\begin{equation}\label{eq:prefix-escape}
 F\notin\Q(t)\quad\Longrightarrow\quad L_p(F)\longrightarrow\infty
 \quad\text{along the good primes}.
\end{equation}

\begin{theorem}[Own-characteristic collapse]\label{thm:ownp}
Let $F\in\Q[[t]]$.  Suppose that there is an integer $h\geq0$ and infinitely
many primes $p$ of good reduction for which
\begin{equation}\label{eq:ownp}
 \sum_{e=0}^{q_p}P_{p,e}(t)\,\overline F_p(t^{p^e})=0,
 \qquad q_p\geq1,
\end{equation}
where $P_{p,e}\in\F_p[t]$, $\deg P_{p,e}\leq h$, and
$P_{p,0}P_{p,q_p}\neq0$.  Then $F$ is constant.
\end{theorem}

\begin{proof}
For $h<n<p$, every term with $e\geq1$ is invisible in the coefficient of
$t^n$: after multiplication by a polynomial of degree at most $h$, the
nonconstant exponents begin at $p$.  Thus the coefficients of $P_{p,0}$ give
a nonzero constant-coefficient recurrence of width at most $h+1$ for an
initial segment whose length tends to infinity with $p$.
Proposition~\ref{prop:prefix-lifting} makes $F$ rational.

Write $F=A/B$ in lowest terms.  After discarding finitely many primes,
degrees and coprimality are preserved.  In characteristic $p$,
$\overline F_p(t^{p^e})=\overline F_p(t)^{p^e}$.  If $F$ has a pole of
order $r$, the last term of \eqref{eq:ownp} has pole order at least
$rp^{q_p}-h$, while all earlier terms have pole order at most
$rp^{q_p-1}$.  For $p-1>h$ the last term has a unique highest pole order and
cannot cancel.  Thus $B$ is constant.  If $A$ is nonconstant of degree $s$,
the same comparison at infinity gives unique top degree
$sp^{q_p}$ versus at most $sp^{q_p-1}+h$.  Therefore $A$ is constant.
\end{proof}

Both nonzero endpoint conditions in \eqref{eq:ownp} are necessary for the
argument.  The theorem closes a particular changing-base route; it does not
assert a general cross-prime Mahler rigidity principle.

For a good prime $p$, define $H_p(F)$ as the minimum of
$\max_{e:P_e\ne0}\deg P_e$ over all $q\geq1$ and polynomials
$P_e\in\F_p[t]$ with $P_0P_q\ne0$ satisfying, in $\F_p[[t]]$,
\[
 \sum_{e=0}^{q}P_e(t)\overline F_p(t^{p^e})=0.
\]
Set $H_p(F)=+\infty$ if no such relation exists.  Then
Theorem~\ref{thm:ownp} has the useful contrapositive
\begin{equation}\label{eq:ownp-degree-escape}
 F\text{ nonconstant}\quad\Longrightarrow\quad
 H_p(F)\longrightarrow+\infty\quad\text{along the good primes}.
\end{equation}

\section{Sharpness and counterexamples}

\subsection{Full own-characteristic rank one, with an exact degree barrier}
Let
\begin{equation}\label{eq:central}
 C(t)=\sum_{n\geq0}\binom{2n}{n}t^n=(1-4t)^{-1/2}.
\end{equation}
For odd $p$, Lucas' theorem gives
\[
 \binom{2(pn+r)}{pn+r}\equiv
 \begin{cases}
 \binom{2r}{r}\binom{2n}{n},&0\leq r\leq(p-1)/2,\\
 0,&(p+1)/2\leq r<p
 \end{cases}\pmod p.
\]
Thus every own-$p$ Cartier section of $C\bmod p$ is a scalar multiple of
$C$ or zero, and iteration gives
\begin{equation}\label{eq:central-full-rank}
 \dim_{\F_p}\operatorname{span}\{\Lambda_wC:w\text{ an own-$p$ word}\}=1.
\end{equation}
At $p=2$, $C\bmod2=1$, so \eqref{eq:central-full-rank} holds for every
prime, not merely at the first Cartier level.

There is nevertheless an exact coefficient-degree cost.  For odd $p$, any
endpoint-nonzero relation
\[
 \sum_{e=0}^{q}P_e(t)C(t^{p^e})=0,
 \qquad P_e\in\F_p[t],\quad \deg P_e\leq h,\quad P_0P_q\ne0,
\]
must satisfy
\begin{equation}\label{eq:central-degree-lower}
 h\geq\frac{p^{q-1}(p-1)}2.
\end{equation}
Indeed, in the quadratic function field $\F_p(t,C)$, put $x=1-4t$.
Then
\[
 C(t^{p^e})=C(t)x^{-(p^e-1)/2}.
\]
For the $x$-adic valuation, one has
$\operatorname{ord}_x(P_q)\leq\deg P_q\leq h$.  The pole contributed by
the top iterate therefore cannot be cancelled by the preceding iterates
unless \eqref{eq:central-degree-lower} holds.  For
$q=1$ equality is attained by
\begin{equation}\label{eq:central-sharp-relation}
 C(t)=(1-4t)^{(p-1)/2}C(t^p)\pmod p.
\end{equation}
Consequently $H_p(C)=(p-1)/2$ for odd $p$.

Yet $C$ is D-finite and nonrational, so it is not $k$-Mahler in
characteristic zero for any fixed $k\geq2$~\cite{bcr}.  In particular,
rank one in a base varying with $p$ cannot replace the all-depth,
fixed-base hypotheses; the necessary degree escapes exactly as predicted by
\eqref{eq:ownp-degree-escape}.  Since $C$ is integral and algebraic, its
reductions modulo every $p^m$ are $p$-automatic under the standard
prime-power automaticity theorem~\cite{profinite}.  Consequently
Corollary~\ref{cor:algebraic} shows that their normalized minimal automaton
state counts tend to infinity.

\subsection{Rational at every finite depth without uniform complexity}
Fix $p$ and set
\begin{equation}\label{eq:hp}
 H_p(t)=\sum_{r\geq1}p^rt^{r^2},\qquad
 F_p(t)=\frac1{1-2t}+H_p(t).
\end{equation}
Modulo $p^m$, the first series is the polynomial obtained by retaining
$1\leq r<m$; hence $F_p\bmod p^m$ is rational for every $m$.
Nevertheless $H_p$, and therefore $F_p$, is non-Mahler in every integer
base.  Indeed its coefficient height is
\[
 h([t^n]H_p)=
 \begin{cases}
 \sqrt n\log p,&n\text{ is a nonzero square},\\
 0,&\text{otherwise},
 \end{cases}
\]
which is sublinear but not $O((\log n)^2)$.  The height-gap theorem of
Adamczewski, Bell and Smertnig~\cite{abs} rules this out for a rational-
coefficient Mahler series.  Subtracting the rational summand preserves
Mahlerity, so $F_p$ is likewise non-Mahler.

This example shows that ``rational modulo every $p^m$'' has no lifting force
without a uniform bound on the rational, Mahler or Cartier complexity.

\section{An adelic one-pole reversal}

We finish with an analytic--arithmetic consequence of having recognized one
fixed-base Mahler equation.  It is a short assembly of the canonical Mahler
denominator and local pole theorem of Adamczewski--Bell--Smertnig, not a new
Mahler-denominator theorem.

\begin{theorem}[One--Mahler adelic pole reversal]
\label{thm:pole-reversal}
Let
\[
 F(t)=\sum_{n\geq0}a_nt^n\in\Q[[t]]
\]
be $k$-Mahler for some fixed $k\geq2$.  For every rational prime $p$, let
$R_p(F)$ be its $p$-adic Taylor radius and assume $R_p(F)\geq1$.  Suppose
the complex germ of $F$ extends meromorphically to
$\mathbb D=\{t\in\C:|t|<1\}$ and that its set of poles in $\mathbb D$ is
exactly $\{\rho\}$, where $0<\rho<1$.  Then
\[
 \tau=\rho^{-1}
\]
is a Pisot number or a Salem number.
\end{theorem}

\begin{proof}
Let $d(t)$ be the normalized $k$-Mahler denominator of $F$ in the sense of
\cite[Definition~3.9]{abs}.  Since $F$ has rational coefficients, its
defining denominator ideal is Galois invariant; uniqueness of the
normalization gives $d\in\Q[t]$.  By definition, for some $r\geq1$ and
$q_i(t)\in\overline\Q[t]$,
\begin{equation}\label{eq:denominator-identity}
 d(t)F(t)=\sum_{i=1}^{r}q_i(t)F(t^{k^i}).
\end{equation}
Both sides are meromorphic on $\mathbb D$, so the identity theorem extends
\eqref{eq:denominator-identity} throughout $\mathbb D$.
If $d(\rho)\ne0$, every term on the right is holomorphic at $t=\rho$:
indeed $0<\rho^{k^i}<\rho$, and $\rho$ is the only pole of $F$ in the
disk.  Equation~\eqref{eq:denominator-identity} would then make $F$
holomorphic at $\rho$, a contradiction.  Thus $d(\rho)=0$ and $\rho$ is
algebraic.

Let $v$ be a nonarchimedean absolute value of $\overline\Q$.  If
$|\rho|_v<1$, Proposition~6.4 of~\cite{abs} forces the $v$-adic Taylor
radius of $F$ to be strictly smaller than $1$, contrary to the hypothesis.
For rational coefficients this is the radius at the underlying rational
prime.  Hence $|\rho|_v\geq1$ at every finite place, equivalently
$|\tau|_v\leq1$, so $\tau$ is an algebraic integer.

Let $\rho'\ne\rho$ be any algebraic conjugate of $\rho$ under an embedding
into $\C$.  Since
$d\in\Q[t]$, one has $d(\rho')=0$.  The proof of
\cite[Proposition~6.4]{abs} gives the stronger local conclusion that a
nonzero denominator root strictly inside the complex unit disk is an actual
pole of the meromorphic continuation.  Thus $|\rho'|<1$ would produce a
second pole in $\mathbb D$.  Consequently every conjugate
$\tau'\ne\tau$ satisfies $|\tau'|\leq1$.

If all these inequalities are strict, $\tau$ is Pisot.  Otherwise let
$P\in\Z[x]$ be the monic irreducible polynomial of $\tau$ and choose a
conjugate $\beta$ on the unit circle.  Then
$\overline\beta=\beta^{-1}$ is also a root of $P$, so $P$ and its
reciprocal polynomial have a common root.  Irreducibility makes $P$
reciprocal up to sign.  Its roots are closed under inversion; since $\tau$
is the only root outside the unit disk, the only roots off the unit circle
are $\tau$ and $\tau^{-1}$.  Thus $\tau$ is Salem.
\end{proof}

\begin{corollary}[Profinite pole interface]\label{cor:pole-interface}
Assume the analytic and finite-place hypotheses of
Theorem~\ref{thm:pole-reversal}.  It is enough to replace the assumption
that $F$ is already Mahler by either of the following, in one fixed base:
\begin{enumerate}[label=\textup{(\roman*)}]
\item a fixed linear box has divergent total precision as in
      Theorem~\ref{thm:mixed-profile};
\item a fixed blocker box has divergent total precision as in
      Theorem~\ref{thm:mixed-blocker}.
\end{enumerate}
Under either condition, $\rho^{-1}$ is Pisot or Salem.
\end{corollary}

\begin{proof}
In case (i), Theorem~\ref{thm:mixed-profile} and
Lemma~\ref{lem:endpoint} make $F$ Mahler in the fixed base.  In case (ii),
Theorem~\ref{thm:mixed-blocker} does so directly.
Theorem~\ref{thm:pole-reversal} then applies.
\end{proof}

\clearpage
\section{Limits and further questions}

The results separate six levels that should not be conflated:
\begin{center}
\begin{tabular}{>{\raggedright\arraybackslash}p{0.33\textwidth}>{\raggedright\arraybackslash}p{0.58\textwidth}}
\toprule
hypothesis & conclusion\\
\midrule
fixed Mahler box, divergent total local precision & exact $k$-Mahler relation\\
bounded blocker box, divergent total local precision & exact $k$-Mahler relation, with no effective output box\\
either criterion in two independent bases & rationality, without D-finiteness\\
D-finite plus either one-base criterion & rationality\\
D-finite and nonrational & reduced $p$-kernel cardinalities tend to infinity\\
one recognized base plus the adelic unique-pole hypotheses & reciprocal pole is Pisot or Salem\\
\bottomrule
\end{tabular}
\end{center}

The present methods do not give an escape rate, a finite positive certificate
for rationality, or automaticity of reductions for all D-finite series.  The
pole reversal does not construct a persistent Mahler window, prove uniqueness
of an inner pole, or handle a boundary point tending to $1$; in particular it
does not prove the Salem--Boyd closure conjecture.  A natural next problem is
to bound the stabilized Smith exponent in terms of arithmetic data of a
D-finite differential equation.  Another is to make the blocker theorem
effective by bounding the Mahler equation produced from $(k,M,D)$.  Both
require input absent from the compactness and Nullstellensatz proofs here.

\enlargethispage{3\baselineskip}
\scriptsize
\begin{thebibliography}{99}

\bibitem{abs}
B.~Adamczewski, J.~P. Bell and D.~Smertnig,
\emph{A height gap theorem for coefficients of Mahler functions},
J. Eur. Math. Soc. (JEMS) \textbf{25} (2023), no.~7, 2525--2571,
\href{https://doi.org/10.4171/JEMS/1244}{doi:10.4171/JEMS/1244}.

\bibitem{abtwobase}
B.~Adamczewski and J.~P. Bell,
\emph{A problem about Mahler functions},
Ann. Sc. Norm. Super. Pisa Cl. Sci. (5) \textbf{17} (2017), no.~4,
1301--1355,
\href{https://doi.org/10.2422/2036-2145.201606_010}{doi:10.2422/2036-2145.201606\_010}.

\bibitem{becker}
P.-G.~Becker,
\emph{$k$-regular power series and Mahler-type functional equations},
J. Number Theory \textbf{49} (1994), no.~3, 269--286,
\href{https://doi.org/10.1006/jnth.1994.1093}{doi:10.1006/jnth.1994.1093}.

\bibitem{bezivin}
J.-P.~B\'ezivin,
\emph{Sur une classe d'\'equations fonctionnelles non lin\'eaires},
Funkcial. Ekvac. \textbf{37} (1994), no.~2, 263--271.

\bibitem{bcr}
J.~P. Bell, M.~Coons and E.~Rowland,
\emph{The rational--transcendental dichotomy of Mahler functions},
J. Integer Seq. \textbf{16} (2013), Article 13.2.10.

\bibitem{eilenberg}
S.~Eilenberg,
\emph{Automata, Languages, and Machines, Vol. A},
Academic Press, New York, 1974.

\bibitem{profinite}
E.~Rowland and R.~Yassawi,
\emph{Profinite automata},
Adv. in Appl. Math. \textbf{85} (2017), 60--83,
\href{https://doi.org/10.1016/j.aam.2016.11.005}{doi:10.1016/j.aam.2016.11.005}.

\bibitem{diagonals}
E.~Rowland and R.~Yassawi,
\emph{Automatic congruences for diagonals of rational functions},
J. Th\'eor. Nombres Bordeaux \textbf{27} (2015), no.~1, 245--288,
\href{https://doi.org/10.5802/jtnb.901}{doi:10.5802/jtnb.901}.

\bibitem{rowlandyassawi}
E.~Rowland and R.~Yassawi,
\emph{Algebraic power series and their automatic complexity modulo prime
powers}, Math. Comp. (2026),
\href{https://doi.org/10.1090/mcom/4227}{doi:10.1090/mcom/4227}.

\end{thebibliography}

\end{document}
